(1) 两能级工作物质 \(H(\omega)=\hbar\omega\sigma_z/2\)，冷、热浴温度分别为 \(T_c<T_h\)，对应频隙 \(0<\omega_c<\omega_h\)。循环按A(冷等容热化)、B(绝热压缩)、C(热等容热化)、D(绝热膨胀)进行。每个等容冲程完全热化，两次有限时间孤立冲程各有对称的能级间转移概率 \(0\le\xi\le1/2\)。求四角点的平均能量与两等容过程热量。
(2) 推导净输出功 \(W_{\rm out}\)、正功条件和热效率 \(\eta\)。验证无跃迁 \(\xi=0\) 时量子Otto效率 \(1-\omega_c/\omega_h\)，并给出不可逆内摩擦造成的精确功损失。
(3) 计算与两个热浴交换的总熵产生 \(\Sigma=-Q_h/T_h-Q_c/T_c\)，用热态的相对熵证明 \(\Sigma\ge0\)。论证正功热机效率不超过Carnot效率，并对 \(\xi=0\) 的正功条件给出等价的温频隙不等式。
(4) 当 \(\xi\) 变化时，严格求净输出功的单调性及失去正功的临界跃迁概率，讨论弱非绝热 \(\xi\ll1\) 的效率修正；对一个数值自洽的工作区间核验能量守恒和耗散功的符号。
